\documentclass[11pt]{article}
\usepackage{amsmath,amssymb,amsthm,hyperref}
\newtheorem{theorem}{Theorem}
\newtheorem{lemma}[theorem]{Lemma}
\newtheorem{corollary}[theorem]{Corollary}
\title{Eventual arithmetic sufficiency on every positive count ray}
\author{Extension of the positive-prefix completion argument}
\date{30 September 2026}
\begin{document}
\maketitle

Fix a positive integer vector $c=(c_1,\ldots,c_n)$. Set
\[
 R(c)=\prod_{\substack{p\le n\text{ prime}\\
              |\{i:p\nmid c_i\}|\ge p}}p.
\]
The empty product is one. The subset-divisibility characterization says
that the vector $tc$ satisfies every necessary factorial divisibility
condition if and only if $R(c)\mid t$.
Indeed, at least $p$ counts not divisible by $p$ supply a violating
$p$-subset. Conversely, if at most $p-1$ counts are $p$-units, every
$k$-subset with $k\ge p$ has product valuation at least
$k-p+1\ge(k-1)/(p-1)\ge v_p(k!)$; for $k<p$ the requirement is empty.
Applying this criterion after multiplication by $t$ proves the claim.

\begin{theorem}\label{main}
For every positive integer vector $c$ there is a finite threshold
$N(c)$ such that, for every integer $t\ge N(c)$, a permutation-fair
family with face-count vector $tc$ exists if and only if $R(c)\mid t$.
More specifically, some positive integer $A$ has both $Ac$ and
$(A+R(c))c$ feasible. Every multiple of $R(c)$ at least
$A(A-R(c))/R(c)$ is consequently a feasible dilation.
\end{theorem}

The conclusion concerns sufficiently large dilations of a fixed profile.
It does not bound the minimum total number of faces sharply, nor imply
that the undilated arithmetic conditions are sufficient.

Use the reduced squarefree algebra and signed-signature theorem from
\texttt{signed\_fair\_signatures.tex}. Put
\[
 H_c=\sum_i c_iX_i,\qquad T_t^{(c)}=\exp(tH_c),\qquad
 D=\operatorname{lcm}(c_1,\ldots,c_n).
\]
The signature $T_t^{(c)}$ is exactly fairness with counts $tc$ whenever
it is represented by a positive word.

\begin{lemma}[Weighted prefix completion]\label{completion}
Every positive word is a literal prefix of a fair positive word with
count vector $Ac$ for some positive integer $A$. More quantitatively,
if $B\ge1$ and the starting word has at most $Bc_i$ occurrences of
letter $i$, one may take
\[
 A=B\prod_{k=2}^n\bigl(D^k+(n!-1)D\bigr).
\]
\end{lemma}
\begin{proof}
Append letters until the count vector is exactly $Bc$. Define weighted
relabelings by
\[
 \psi_\sigma(X_i)=\frac{c_{\sigma(i)}}{c_i}X_{\sigma(i)}
 \qquad(\sigma\in S_n).
\]
These are a group action by algebra automorphisms, they fix $H_c$, and
they are conjugate to ordinary relabelings by the diagonal change
$X_i\mapsto c_iX_i$. Thus a homogeneous Lie polynomial of degree
$k>1$ invariant under every $\psi_\sigma$ is zero, by the invariant-Lie
lemma in the positive-completion note.

Suppose the current positive word $w$ has
\[
 \log\operatorname{sig}(w)=BH_c+E_k+\text{terms of degree greater than }k,
\]
with all errors in degrees $2,\ldots,k-1$ zero.
Let $\delta_D$ be the dilation $X_i\mapsto DX_i$.
For every nonidentity permutation $\sigma$, the image of $w$ under
$\delta_D\psi_\sigma$ is an ordinary positive word: replace each
letter $i$ by $D c_{\sigma(i)}/c_i$ copies of $\sigma(i)$.
All these multiplicities are positive integers.

Concatenate first $D^k$ copies of $w$, and then these transformed
words for every $\sigma\ne\mathrm{id}$. The first block has logarithm
\[
 D^k BH_c+D^k E_k+\cdots,
\]
whereas the block for $\sigma\ne\mathrm{id}$ has logarithm
\[
 DBH_c+D^k\psi_\sigma(E_k)+\cdots.
\]
All degree-one terms are proportional to $H_c$ and commute.
Any BCH bracket involving an error has degree at least $k+1$.
The new degree-$k$ error is therefore
$D^k\sum_{\sigma\in S_n}\psi_\sigma(E_k)=0$.
The new first-degree multiplier is
$B(D^k+(n!-1)D)$, and the old word remains a literal prefix.
Iterating for $k=2,\ldots,n$ gives the claim.
\end{proof}

\begin{lemma}[Clearing signed words on a prescribed ray]\label{clearing}
For every signed group element $g$ there is a positive integer $A$
such that both $T_A^{(c)}$ and $gT_A^{(c)}$ are represented by positive
words.
\end{lemma}
\begin{proof}
Weighted completion supplies, for any positive signatures $u,v$,
positive complements $u',v'$ with $uu'=T_a^{(c)}$ and
$vv'=T_b^{(c)}$. These fair signatures commute, so
\[
 u(u'T_b^{(c)})=v(v'T_a^{(c)})=T_{a+b}^{(c)}.
\]
Hence positive signatures have common right multiples. As in the
equal-count case, every signed element is a right fraction $xy^{-1}$
of positive signatures. Complete $yz=T_A^{(c)}$; then
$gT_A^{(c)}=xz$ is positive.
\end{proof}

\begin{proof}[Proof of Theorem~\ref{main}]
For a prime $p$, multiplication of all counts by a dilation $t$ either
makes every count divisible by $p$ (if $p\mid t$), or leaves unchanged
the set of $p$-unit counts. The subset-divisibility characterization
therefore gives exactly $R(c)\mid t$. In particular the signed-signature
theorem gives $T_{R(c)}^{(c)}$ as a signed group element.
Apply Lemma~\ref{clearing} to this element. The resulting positive
signatures are $T_A^{(c)}$ and $T_{A+R(c)}^{(c)}$.
Necessity gives $R(c)\mid A$.

Concatenations of these two fair words remain on the same fair ray.
Writing $a=A/R(c)$, they realize every nonnegative integer combination
of $a$ and $a+1$, multiplied by $R(c)$. Every integer at least
$a(a-1)$ is such a combination, proving the stated threshold.
\end{proof}

\paragraph{Review status.}
The full weighted-power proof has an independent line-by-line
mathematical review, and exact reduced-algebra verification has passed
on six small count rays; see
\texttt{../positive\_saturation/verify\_weighted\_completion.py}.
No machine-checked formal verification is claimed.
\end{document}
